\section*{Bab \ref{ch:g12:early-earth} --- Bumi Purba dan Atmosfernya}

\begin{solution}{exo:g12:early-earth:1}
(a) \omterm{def:g12:early-earth:primitive-atmosphere}{pelepasan gas}; (b) \omterm{def:g12:early-earth:stromatolite}{stromatolit}; (c) \omterm{def:g12:early-earth:bif}{formasi besi berpita}; (d) \omterm{def:g12:early-earth:red-beds}{lapisan
merah}; (e) \omterm{def:g12:early-earth:goe}{Peristiwa Oksidasi Besar}.
\end{solution}

\begin{solution}{exo:g12:early-earth:2}
Zirkon yang menunjukkan air cair (4.4 miliar tahun); \omterm{def:g12:early-earth:stromatolite}{stromatolit} pertama di
Pilbara (3.43); \omterm{def:g12:early-earth:bif}{formasi besi berpita} Hamersley (lebih dari 2.5); berakhirnya
\omterm{def:g12:early-earth:goe}{Peristiwa Oksidasi Besar} (2.06); berakhirnya sebagian besar \omterm{def:g12:early-earth:bif}{formasi besi
berpita} (1.8).
\end{solution}

\begin{solution}{exo:g12:early-earth:3}
Pada $t = 0$, $L = L_\odot/(1 + 0.4) = L_\odot/1.4 \approx 0.714\,L_\odot$:
sekitar \qty{29}{\percent} lebih redup, hampir sepertiga kurang terang
daripada sekarang.
\end{solution}

\begin{solution}{exo:g12:early-earth:4}
Bahwa udara waktu itu hampir tidak mengandung oksigen. Pirit dihancurkan oleh
oksigen; butir-butirnya hanya dapat digulirkan sungai dan diendapkan dalam
keadaan utuh bila udara dan air di permukaan tidak mengandung oksigen.
\end{solution}

\begin{solution}{exo:g12:early-earth:5}
Uap air, yang mengembun menjadi samudra, dan karbon dioksida, yang tinggal di
udara dalam jumlah yang jauh lebih besar daripada sekarang dan menghangatkan
Bumi melalui efek rumah kacanya.
\end{solution}

\begin{solution}{exo:g12:early-earth:6}
3 miliar tahun yang lalu: $1 - t/t_\odot = 3/4.567 \approx 0.657$, $L \approx
1/1.263 \approx 0.792\,L_\odot$, $T \approx 255 \times 0.792^{1/4} \approx
\qty{241}{K}$. 2.45 miliar tahun yang lalu: $1 - t/t_\odot \approx 0.536$, $L
\approx 1/1.215 \approx 0.823\,L_\odot$, $T \approx 255 \times 0.823^{1/4}
\approx \qty{243}{K}$.
\end{solution}

\begin{solution}{exo:g12:early-earth:7}
\qty{1}{kg} besi adalah $1000/55.85 \approx 17.9$ mol; satu molekul oksigen
mengikat empat ion besi, sehingga $17.9/4 \approx 4.48$ mol \ce{O2}, yaitu
$4.48 \times 32.00 \approx \qty{143}{g}$ oksigen.
\end{solution}

\begin{solution}{exo:g12:early-earth:8}
(1) $2.46 - 2.06 = 0.40$ miliar tahun, 400 juta tahun. (2) Dari 4.54 sampai
2.46 miliar tahun yang lalu: $2.08/4.54 \approx \qty{46}{\percent}$ sejarah
Bumi. (3) Sekitar $3.43 - 2.46 \approx 1$ miliar tahun: oksigen yang dibuat
mikroba terpakai habis untuk waktu yang lama sebelum dapat tinggal di udara.
\end{solution}

\begin{solution}{exo:g12:early-earth:9}
\omterm{def:g12:early-earth:bif}{Formasi besi berpita} menunjukkan laut dalam tanpa oksigen, yang dapat membawa
besi terlarut, di bawah laut dangkal yang diterangi Matahari dengan oksigen
yang cukup untuk mengendapkannya. \omterm{def:g12:early-earth:red-beds}{Lapisan merah}, yang diendapkan di darat,
menunjukkan oksigen di udara dan di air permukaan, cukup untuk membuat besi
pada butir-butirnya berkarat.
\end{solution}

\begin{solution}{exo:g12:early-earth:10}
Di Hamelin Pool, Teluk Shark, Australia Barat, yang airnya dua kali lebih asin
daripada air laut biasa. Hamparan mikroba, terutama sianobakteri, menjebak dan
mengikat butir-butir sedimen dan menumbuhkan karbonat; hamparan baru tumbuh di
atasnya, dan seterusnya, lapis demi lapis.
\end{solution}

\begin{solution}{exo:g12:early-earth:11}
(1) $21 \times 10^{-5} = \qty{0.00021}{\percent}$. (2) $21/10 =
\qty{2.1}{\percent}$.
\end{solution}

\begin{solution}{exo:g12:early-earth:12}
\omterm{def:g10:habitable-planet:equilibrium-temperature}{Suhu kesetimbangan} tidak berubah bila $L(1 - A')$ sama dengan $1 - A$:
$1 - A' = 0.706/0.741 \approx 0.953$, sehingga $A' \approx 0.047$. Bumi harus
memantulkan kurang dari \qty{5}{\percent} sinar Matahari, hampir seperti benda
hitam, hal yang tidak mungkin dilakukan sebuah planet bersamudra yang mempunyai
awan dan sedikit es. Jalan keluar yang lain adalah efek rumah kaca yang lebih
kuat.
\end{solution}

\begin{solution}{exo:g12:early-earth:13}
(1) $273 - 237 = \qty{36}{K}$ untuk mencapai titik beku; $288 - 237 = \qty{51}{K}$
untuk mencapai rata-rata saat ini. (2) Efek rumah kaca saat ini menambahkan $288 - 255 =
\qty{33}{K}$: Bumi muda memerlukan efek yang lebih kuat. (3) \qty{0.1}{bar}
adalah sepersepuluh tekanan seluruh atmosfer saat ini, untuk karbon dioksida
saja.
\end{solution}

\begin{solution}{exo:g12:early-earth:14}
(1) Sebelum 2450 juta tahun yang lalu tanda itu ada, sehingga udara mengandung
oksigen kurang dari $10^{-5}$ kadar saat ini; setelah 2090 juta tahun tanda itu
tidak ada, sehingga udara mengandung lebih banyak. (2) Ya: \omterm{def:g12:early-earth:goe}{Peristiwa Oksidasi
Besar} berlangsung dari sekitar 2460 sampai 2060 juta tahun yang lalu, selang
yang sama.
\end{solution}

\begin{solution}{exo:g12:early-earth:15}
(1) Sekitar $4567 - 60 = 4507$ juta tahun yang lalu. (2) Paling lama $4507 - 4404 =
103$ juta tahun setelah tumbukan: Bumi sudah mempunyai kerak dan air dalam
waktu sekitar seratus juta tahun setelah tumbukan itu.
\end{solution}

\begin{solution}{pb:g12:early-earth:1}
\textbf{1.} $4.567 - 4.0 = 0.567$ miliar tahun.

\textbf{2.} $1 - 0.567/4.567 \approx 1 - 0.124 = 0.876$.

\textbf{3.} $L = 1/(1 + 0.4 \times 0.876) = 1/1.350 \approx 0.741$ dari
luminositas sekarang.

\textbf{4.} $0.741 \times 1361 \approx \qty{1008}{W/m^2}$.

\textbf{5.} $1 - 0.741 \approx 0.26$: \qty{26}{\percent} lebih redup.

\textbf{6.} $T_{\mathrm{eq}} = \left(S(1-A)/(4\sigma)\right)^{1/4}$.

\textbf{7.} Dengan $A$ dan $\sigma$ tetap, $T_{\mathrm{eq}}$ sama dengan sebuah
konstanta dikalikan $S^{1/4}$, sehingga $T'/T = (S'/S)^{1/4}$.

\textbf{8.} $0.741^{1/4} = \sqrt{\sqrt{0.741}} \approx \sqrt{0.861} \approx
0.928$.

\textbf{9.} $255 \times 0.928 \approx \qty{237}{K}$.

\textbf{10.} $237 - 273 \approx \qty{-36}{\celsius}$.

\textbf{11.} Bahwa ada air cair di permukaan, sebuah samudra, 4.4 miliar tahun
yang lalu.

\textbf{12.} Di air dangkal, seperti \omterm{def:g12:early-earth:stromatolite}{stromatolit} hidup saat ini.

\textbf{13.} \omterm{def:g12:early-earth:bif}{Formasi besi berpita} Pegunungan Hamersley, yang diendapkan di
dasar laut.

\textbf{14.} Ya: \qty{237}{K} jauh di bawah \qty{273}{K}; dengan atmosfer
seperti sekarang permukaan rata-ratanya akan membeku.

\textbf{15.} Batuan menunjukkan air cair di bawah Matahari yang terlalu redup
untuk menjaganya tetap cair dengan atmosfer seperti sekarang: \omterm{def:g12:early-earth:faint-young-sun}{masalah Matahari
muda yang redup}.

\textbf{16.} $288 - 255 = \qty{33}{K}$.

\textbf{17.} Setidaknya $273 - 237 = \qty{36}{K}$, lebih banyak daripada saat
ini.

\textbf{18.} $A' = 1 - 0.706/0.741 \approx 0.047$: planet yang hampir hitam,
tidak masuk akal.

\textbf{19.} Karbon dioksida, yang jauh lebih melimpah daripada sekarang;
metana, yang dibuat oleh mikroba awal, mungkin ikut membantu.

\textbf{20.} \omterm{def:g10:habitable-planet:equilibrium-temperature}{Suhu kesetimbangan} Bumi empat miliar tahun yang lalu adalah
\textbf{sekitar 237 K}, kira-kira \qty{36}{\celsius} di bawah titik beku.
\end{solution}
